\section{本章小结}
\label{sec:opt-mmb-summary}

章首提出的问题是：当候选的好坏不再由一个数直接决定时，“最优答案”是什么，怎样算出它，又怎样检查它。本章对三种结构分别作了回答。

多个指标共存时，支配关系先排除全面劣势的候选，剩下的Pareto集合通常不止一点。非空紧可行域与连续指标保证正权加权和有最小点，从而保证Pareto点存在。加权和用价格表达偏好，只能到达可达集合左下凸包上的点，凹陷中的Pareto点永远选不到，权重变化时解还可能跳变；阈值约束用上限表达偏好，能到达每个Pareto点，但不唯一的解只保弱Pareto性。不预设偏好时，梯度凸包中的最小范数元给出共同下降方向，它为零只说明一阶意义下不能同时改善，凸性与正权才把它加强为Pareto最优。

两组变量方向相反时，先规定次序：外层最小化的是对手最优响应构成的上包络，内层近似求解只给出它的下界。鞍点表示双方互为全局最优响应，间隙为零是它的可计算证书；紧凸凹条件保证上下值相等且鞍点存在，但不保证任何特定算法能找到它。双线性例中，同时梯度下降—上升向外螺旋，加上投影后循环，交替版本有界打转；外梯度法先试探、再从旧点校正，在凸凹且Lipschitz的条件下平均点的间隙以$O(1/T)$下降。

一个问题嵌套另一个问题的解时，先规定内层怎样选解：唯一解、乐观或悲观规则、指定算法的输出，各自给出不同的外层目标。迭代展开对有限步目标$J_T$精确求导，代价是保存或重算轨迹；隐式微分对局部分支目标$\widehat J$求导，只需解点处的伴随方程，但要求Hessian可逆，误差同时来自内层残差与伴随残差，小曲率会放大二者。展开与隐式两种导数在正则参数例中分别为$-15/16$与$-6$，二者都正确，差别在于目标不同。外层求解可以直接调用前面章节的光滑、投影和随机方法，但超梯度误差限定了外层能达到的精度，而把局部分支误当作全局选解的错误，任何计算精度都无法弥补。

三类问题共同提醒我们分开陈述四件事：答案是否存在，局部判据是否满足，算法是否在进展，以及有没有可以核对的全局证书。前面各章的标量方法在这里分别承担了一部分工作：加权和与阈值问题直接交给光滑与约束方法求解，投影与单调性支撑了外梯度的分析，二次上界与隐函数定理支撑了超梯度的计算与外层下降。下一章转向另一种扩展：未知量从有限维参数变成由数据决定的函数，双层结构也会在正则参数的选择中再次出现。
